\documentclass[12pt,a4paper]{report}
\usepackage{fontspec}
\setmainfont{Times New Roman}

\usepackage[vietnamese]{babel}
\shorthandoff{"}
\usepackage{amsmath,amsfonts,amssymb}
\usepackage{geometry}
\geometry{left=3cm,right=2cm,top=2cm,bottom=2cm}
\usepackage{setspace}
\onehalfspacing

\usepackage{tabularx}
\usepackage{longtable}
\usepackage{booktabs}
\usepackage{colortbl}
\usepackage{array}
\usepackage{hyperref}
\hypersetup{hidelinks}
\usepackage{url}
\usepackage{titlesec}

\titleformat{\chapter}[display]
  {\normalfont\bfseries\Large}{\chaptertitlename\ \thechapter}{10pt}{\Large}
\titlespacing*{\chapter}{0pt}{0pt}{20pt}

\begin{document}
\setcounter{chapter}{4}
\chapter{XÂY DỰNG CHƯƠNG TRÌNH MINH HỌA}

\section{Thư viện và công cụ sử dụng}

Trong quá trình phát triển ứng dụng quản lý lớp học thêm (hệ thống ``Lớp Mình''), nhóm chúng em đã sử dụng các ngôn ngữ lập trình, framework, thư viện, công cụ phát triển và cơ sở dữ liệu phù hợp để đảm bảo tính hiệu quả, bảo mật và dễ bảo trì. Các công cụ được chọn dựa trên yêu cầu dự án, bao gồm frontend xây dựng trên nền tảng Web hiện đại (HTML5, CSS3, JavaScript ES6+, Feather Icons), backend dựa trên Python với framework FastAPI và máy chủ Uvicorn, cùng hệ quản trị cơ sở dữ liệu quan hệ mạnh mẽ PostgreSQL (v14+). Dưới đây là danh sách chi tiết các thư viện và công cụ được sử dụng, được trình bày dưới dạng bảng.

\setlength{\tabcolsep}{3pt}
\renewcommand{\arraystretch}{1.25}
\begin{longtable}{|>{\raggedright\arraybackslash}p{3.8cm}|>{\raggedright\arraybackslash}p{5.4cm}|>{\raggedright\arraybackslash}p{6.4cm}|}
    \caption{Thư viện và công cụ sử dụng}\label{tab:c5-thu-vien-cong-cu}\\
    \hline\rowcolor{blue!11}
    \multicolumn{1}{|c|}{\textbf{Mục đích}} & \multicolumn{1}{c|}{\textbf{Công cụ}} & \multicolumn{1}{c|}{\textbf{Địa chỉ URL}} \\ \hline
    \endfirsthead

    \multicolumn{3}{r}{\textit{Tiếp bảng 5.1: Thư viện và công cụ sử dụng}}\\
    \hline\rowcolor{blue!11}
    \multicolumn{1}{|c|}{\textbf{Mục đích}} & \multicolumn{1}{c|}{\textbf{Công cụ}} & \multicolumn{1}{c|}{\textbf{Địa chỉ URL}} \\ \hline
    \endhead

    \hline
    \multicolumn{3}{r}{\footnotesize\textit{Còn tiếp trang sau}}\\
    \endfoot

    \hline
    \endlastfoot

    Môi trường & Python (v3.11+) \newline Node.js (v20+) & \url{https://www.python.org/} \newline \url{https://nodejs.org/} \\ \hline
    Ngôn ngữ lập trình & Python (v3.11.6) \newline JavaScript (ES6+ / CommonJS) \newline HTML5 / CSS3 \newline SQL (PostgreSQL Dialect) & \url{https://www.python.org/} \newline \url{https://developer.mozilla.org/en-US/docs/Web/JavaScript} \newline \url{https://developer.mozilla.org/} \newline \url{https://www.postgresql.org/docs/} \\ \hline
    Framework frontend & HTML5 Semantic \& Vanilla CSS3 \newline Hệ thống CSS Variables, Flexbox, Grid & \url{https://developer.mozilla.org/} \newline \url{https://developer.mozilla.org/en-US/docs/Web/CSS} \\ \hline
    Build tool \& Server frontend & Python HTTP Server (http.server v3.11+) & \url{https://docs.python.org/3/library/http.server.html} \\ \hline
    CSS framework & Custom Modular CSS Architecture (variables, base, components, dashboard) & \url{https://developer.mozilla.org/en-US/docs/Web/CSS} \\ \hline
    UI icons library & Feather Icons (v4.29.0) & \url{https://feathericons.com/} \\ \hline
    Charts \& Visualization & HTML5 Canvas \& Dynamic SVG & \url{https://developer.mozilla.org/en-US/docs/Web/API/Canvas_API} \\ \hline
    Routing \& Access control & Client-side Router \& Auth Guard (auth.js) & \url{https://developer.mozilla.org/en-US/docs/Web/API/History_API} \\ \hline
    Notification library & Custom Modal \& Toast Notifications (notifications.js) & \url{https://developer.mozilla.org/en-US/docs/Web/API/Notifications_API} \\ \hline
    HTTP client & Fetch API (Native Browser Web API) & \url{https://developer.mozilla.org/en-US/docs/Web/API/Fetch_API} \\ \hline
    Typography / Phông chữ & Google Fonts (Inter font family) & \url{https://fonts.google.com/specimen/Inter} \\ \hline
    Image upload \& File handling & FormData API \& File API (Native Web API) & \url{https://developer.mozilla.org/en-US/docs/Web/API/FormData} \\ \hline
    Client-side storage & Web Storage API (LocalStorage) & \url{https://developer.mozilla.org/en-US/docs/Web/API/Window/localStorage} \\ \hline
    Environment variables & Python os \& Dotenv configuration (.env) & \url{https://fastapi.tiangolo.com/advanced/settings/} \\ \hline
    Framework backend & FastAPI (v0.115.0+) & \url{https://fastapi.tiangolo.com/} \\ \hline
    Web server backend & Uvicorn [standard] (v0.30.0+) & \url{https://www.uvicorn.org/} \\ \hline
    Trình kết nối CSDL & Psycopg 3 (psycopg[binary] v3.2.0+) & \url{https://www.psycopg.org/} \\ \hline
    ORM \& Data Validation & Pydantic (v2.x) & \url{https://docs.pydantic.dev/} \\ \hline
    Authentication \& Security & Hashlib \& Secrets (scrypt KDF, Salt 16-byte, SHA-256) & \url{https://docs.python.org/3/library/hashlib.html} \\ \hline
    Session management & Bearer Session Token (SHA-256 token hash) & \url{https://datatracker.ietf.org/doc/html/rfc6750} \\ \hline
    File upload backend & Python-multipart (v0.0.12+) & \url{https://github.com/Kludex/python-multipart} \\ \hline
    Hỗ trợ \& Liên hệ trung tâm & Hotline 028 3822 8888 \& Trung tâm hỗ trợ 3 vai trò & \url{https://lopminh.edu.vn/} \\ \hline
    Caching \& Session & PostgreSQL Sessions (auth\_sessions) \& LocalStorage & \url{https://www.postgresql.org/} \\ \hline
    Unique ID generator & UUID v4 (Python uuid module) & \url{https://docs.python.org/3/library/uuid.html} \\ \hline
    CORS handling & FastAPI CORSMiddleware & \url{https://fastapi.tiangolo.com/tutorial/cors/} \\ \hline
    Rate limiting \& Brute-force & Lockout sau 3 lần sai / 15 phút (FastAPI \& Database) & \url{https://fastapi.tiangolo.com/} \\ \hline
    API documentation & Swagger UI / OpenAPI 3.0 \newline ReDoc & \url{https://swagger.io/} \newline \url{https://redocly.com/} \\ \hline
    Development server & Uvicorn Auto-reload (Backend) \newline Python HTTP Server (Frontend) & \url{https://www.uvicorn.org/} \newline \url{https://docs.python.org/3/library/http.server.html} \\ \hline
    Kiểm thử tự động & Python Unittest (API contract, DB schema, security) \newline Node.js Test Runner (node --test) & \url{https://docs.python.org/3/library/unittest.html} \newline \url{https://nodejs.org/api/test.html} \\ \hline
    Cơ sở dữ liệu chính & PostgreSQL (v14+) & \url{https://www.postgresql.org/} \\ \hline
    CSDL cấu trúc linh hoạt & PostgreSQL JSONB & \url{https://www.postgresql.org/docs/current/datatype-json.html} \\ \hline
    Kịch bản điều phối & PowerShell (run\_project.ps1) & \url{https://learn.microsoft.com/powershell/} \\ \hline
    IDE lập trình & Visual Studio Code (v1.90+) & \url{https://code.visualstudio.com/} \\ \hline
    Version control & Git (v2.56.0.windows.4) & \url{https://git-scm.com/} \\ \hline
    Remote repository & GitHub & \url{https://github.com/} \\ \hline
    Package manager & PIP (Python Package Installer v23.2+) & \url{https://pip.pypa.io/} \\ \hline
    Công cụ soạn thảo báo cáo & XeLaTeX / LaTeX (TeX Live) \newline Microsoft Word & \url{https://www.latex-project.org/} \newline \url{https://www.microsoft.com/microsoft-365/word} \\ \hline
\end{longtable}

\section{Kết quả chương trình minh họa}

Ứng dụng quản lý lớp học thêm (``Lớp Mình'') đã được phát triển thành một hệ thống web hoàn chỉnh, mang lại giải pháp toàn diện cho việc tổ chức, điều hành và theo dõi các hoạt động dạy và học tại các lớp học tư thục từ đầu đến cuối, với khả năng hỗ trợ đa người dùng qua ba vai trò chuyên biệt: Giáo viên/Gia sư, Học sinh và Phụ huynh.

Kết quả đạt được là một sản phẩm phần mềm có cấu trúc module hóa rõ ràng, đóng gói hoàn chỉnh và sẵn sàng vận hành. Toàn bộ mã nguồn frontend được xây dựng trên nền tảng web hiện đại với HTML5, CSS3 có hệ thống biến màu sắc và bố cục nhất quán, kết hợp JavaScript ES6+ và Feather Icons nhằm mang lại trải nghiệm tương tác mượt mà, tối ưu hiệu năng tải trang và dễ dàng bảo trì mà không phụ thuộc vào framework cồng kềnh. Backend được phát triển trên nền Python với FastAPI và máy chủ Uvicorn, triển khai đầy đủ các endpoint RESTful an toàn, có tài liệu trực quan tự động qua Swagger UI. Cơ sở dữ liệu quan hệ PostgreSQL (v14+) được chuẩn hóa với các ràng buộc toàn vẹn, chỉ mục tối ưu, view tự động đối soát công nợ học phí (\texttt{v\_tuition\_invoice\_balances}), cùng các cơ chế bảo mật nghiêm ngặt như băm mật khẩu bằng thuật toán scrypt với salt ngẫu nhiên 16 byte, cơ chế tự động khóa đăng nhập 15 phút sau 3 lần thất bại liên tiếp, và quản lý phiên bằng Bearer session token băm SHA-256 lưu trong cơ sở dữ liệu. Toàn bộ hệ thống được bảo đảm chất lượng bởi kịch bản tự động hóa PowerShell và bộ kiểm thử tự động toàn diện kết hợp giữa Python \texttt{unittest} và Node.js test runner.

Ý nghĩa của sản phẩm nằm ở việc số hóa và chuẩn hóa toàn diện quy trình quản lý lớp học thêm truyền thống, xóa bỏ hoàn toàn các sai sót thủ công trong ghi chép sổ sách và tính toán học phí, đồng thời tăng cường sự kết nối mật thiết giữa giáo viên, phụ huynh và học sinh thông qua cơ chế thông báo và tra cứu tức thời. Vai trò cốt lõi của ứng dụng là cung cấp một nền tảng vận hành tập trung, minh bạch: hỗ trợ giáo viên tinh gọn tối đa các tác vụ hành chính (điểm danh một chạm, quản lý ngân hàng bài tập, sổ điểm điện tử, lập hóa đơn tự động); mang lại cho phụ huynh cổng thông tin đáng tin cậy để giám sát sát sao chuyên cần, tiến độ học tập và tình trạng học phí của con em; đồng thời trao quyền cho học sinh chủ động nắm bắt lịch học, nộp bài tập và theo dõi biểu đồ tiến bộ học tập của chính mình.

\end{document}
